% ============================================================
%  results.tex  —  §7
%  Target venue: Omega
%  Organisation: mechanism first, benchmark as confirmation
%  Updated: 29/Sep/2026
% ============================================================

\section{Results}
\label{sec:results}

% ----------------------------------------------------------
\subsection{Base portfolios under nominal weights}
\label{sec:base}
% ----------------------------------------------------------

Table~\ref{tab:canonical} reports the portfolio selected by each
method--assembly pair under the nominal AHP weights $\wahp$, with
normalisation constants computed over all $\Ncand$ candidates (space~$A$).

The reproduction is exact: SAW under optimal assembly returns
$\{184,430,460\}$, the best-compromise portfolio of the source study,
with objective values agreeing to $10^{-16}$.
The pipeline therefore recovers the published solution before any
factor is varied.

Two failures are visible at nominal weights.
Greedy assembly of the \emph{same} SAW score returns $\{184,279,320\}$
with $\mathrm{gap}=0.032$, a dominated portfolio sharing one site with
the optimum.
TOPSIS with vector normalisation and VIKOR are also dominated
($\mathrm{gap}=0.048$ and $0.052$), while TOPSIS with min--max
normalisation and PROMETHEE~II attain the reference set.

\begin{table*}[t]
\centering
\caption{Portfolios under nominal weights $\wahp$, space~$A$
  (normalisation over all $\Ncand$ candidates). The gap is defined
  in~\eqref{eq:gap}; $\mathrm{gap}=0$ iff the portfolio lies in the
  reference set.}
\label{tab:canonical}
\begin{tabular*}{\textwidth}{@{\extracolsep{\fill}}llccccc}
\toprule
Method & Assembly & Sites & $F_1$ & $F_2$ & $F_3$ & gap \\
\midrule
SAW, min--max     & greedy  & $\{184,279,320\}$ & 0.1273 & 0.5963 & 0.6207 & 0.032 \\
SAW, min--max     & optimal & $\{184,430,460\}$ & 0.0811 & 0.5568 & 0.6170 & 0.000 \\
TOPSIS, vector    & greedy  & $\{233,320,355\}$ & 0.1113 & 0.5669 & 0.5358 & 0.048 \\
TOPSIS, min--max  & greedy  & $\{184,227,460\}$ & 0.1104 & 0.6712 & 0.6883 & 0.000 \\
VIKOR             & greedy  & $\{233,279,320\}$ & 0.1190 & 0.5819 & 0.5879 & 0.052 \\
PROMETHEE~II      & greedy  & $\{51,233,345\}$  & 0.0447 & 0.5900 & 0.4852 & 0.000 \\
\midrule
Source study      & ---     & $\{184,430,460\}$ & 0.0811 & 0.5568 & 0.6170 & 0.000 \\
\bottomrule
\end{tabular*}
\end{table*}

A first indication that the reference set is itself implicated: when
normalisation constants are recomputed over the $\Nvia$ admissible sites
(space~$A'$), \emph{all six} portfolios become non-dominated.
The failures in space~$A$ are not solely attributable to the methods;
they are partly induced by computing the value scale over sites that
cannot be selected.
Section~\ref{sec:refset} returns to this.

% ----------------------------------------------------------
\subsection{Confirmation of the amplitude lemma}
\label{sec:lemmaconf}
% ----------------------------------------------------------

Lemma~\ref{lem:amplitude} predicts that performance degrades with the
distortion the normalisation imposes on the elicited weights, measured
on the simplex by
\begin{equation}
  \label{eq:TV_results}
  \mathrm{TV} = \tfrac12\sum_{j=1}^{m}\bigl|w_j^{\mathrm{eff}}-w_j\bigr|,
  \qquad
  w_j^{\mathrm{eff}} = \frac{w_j A_j}{\sum_k w_k A_k}.
\end{equation}
Total variation is the appropriate measure because, as noted after
Corollary~\ref{cor:three}, $w^{\mathrm{eff}}$ is invariant to a common
rescaling of all amplitudes: a factor $c>0$ applied to every $A_j$
cancels in~\eqref{eq:TV_results} and produces no re-weighting.

Table~\ref{tab:tvorder} pairs the predicted distortion with observed
performance, restricted to additive aggregation --- the scope in which
the lemma is exact.

\begin{table}[t]
\centering
\caption{Predicted distortion and observed performance under additive
  aggregation (SAW), $\kappa=200$, space~$A'$.
  Amplitudes $A_j$ are computed over the $\Nvia$ admissible sites.}
\label{tab:tvorder}
\begin{tabular}{lccc}
\toprule
Normalisation & $A_j$ & $\mathrm{TV}$ & $Q_2[A']$ \\
\midrule
min--max & $\equiv 1$                        & 0.000 & 82.3\% \\
midranks & $(r_j^{\max}-r_j^{\min})/(N{+}1)$ & 0.091 & 76.5\% \\
vector   & $\rho_j$                          & 0.312 & 55.5\% \\
\bottomrule
\end{tabular}
\end{table}

The ordering is strictly monotone in the predicted direction, and the
magnitudes are approximately proportional: a distortion of $0.091$
costs $5.8$~pp, while $0.312$ costs $26.8$~pp.
A least-squares fit over the three levels gives a slope of
$-84.1$~pp per unit of total variation with $R^2=0.995$.
We report the \emph{ordering and approximate proportionality}, not a
predictive coefficient: with three levels the residual degree of freedom
is one, and any monotone triple admits a close linear fit.
The claim we make is the analytic one of Corollary~\ref{cor:three} ---
min--max is the unique weight-preserving choice among the three ---
for which the empirical ordering is corroboration rather than
evidential foundation.

Two entries of Table~\ref{tab:tvorder} deserve comment because they
isolate distinct mechanisms acting on the same criterion.
Under vector normalisation, \texttt{perm} is \emph{inflated} to
$2.865\times$ its nominal weight, because its long right tail
concentrates the Euclidean norm in few sites.
Under midranks scaled by $N+1$, the same criterion is \emph{deflated}
to $\approx0.55\times$, because its 350 tied zeros compress the
attainable rank range to $[0.003,\,0.550]$ while every other criterion
spans $[0.003,\,0.997]$.
The two normalisations distort the same criterion in opposite
directions, for reasons the lemma makes explicit and neither
normalisation makes visible to the analyst.

% ----------------------------------------------------------
\subsection{The factorial design: aggregation is downstream}
\label{sec:factorial}
% ----------------------------------------------------------

Table~\ref{tab:factorial} reports the balanced $2^3$ design and
Table~\ref{tab:effects} the contrasts, with $\nboot$-replicate bootstrap
intervals paired by draw.

\begin{table}[t]
\centering
\caption{$Q_2[A']$ for the balanced $2^3$ design ($\kappa=200$,
  $\ndraws$ draws). The SAW/min--max/optimal cell is fixed at $100\%$ by
  Proposition~\ref{prop:p1} and serves as a regression test.}
\label{tab:factorial}
\begin{tabular}{llc}
\toprule
Aggregation & Normalisation, assembly & $Q_2[A']$ \\
\midrule
\multirow{4}{*}{SAW}
 & min--max, greedy  & 64.6\% \\
 & min--max, optimal & \textbf{100.0\%} \\
 & vector, greedy    & 54.2\% \\
 & vector, optimal   & 56.8\% \\
\midrule
\multirow{4}{*}{TOPSIS}
 & min--max, greedy  & 72.4\% \\
 & min--max, optimal & 96.6\% \\
 & vector, greedy    & 50.6\% \\
 & vector, optimal   & 52.4\% \\
\bottomrule
\end{tabular}
\end{table}

\begin{table}[t]
\centering
\caption{Contrasts in percentage points of $Q_2[A']$ ($\kappa=200$;
  $\nboot$ bootstrap replicates paired by draw).}
\label{tab:effects}
\begin{tabular}{lrc}
\toprule
Contrast & Effect & 95\% CI \\
\midrule
Normalisation (min--max $-$ vector) & $+29.9$ & $[+27.1,\,+32.9]$ \\
Assembly (optimal $-$ greedy)       & $+16.0$ & $[+14.1,\,+17.9]$ \\
Aggregation (SAW $-$ TOPSIS)        & $+0.9$  & $[-0.5,\,+2.3]$ \\
\midrule
Assembly within min--max            & $+29.8$ & $[+27.4,\,+32.2]$ \\
Assembly within vector              & $+2.2$  & $[-0.5,\,+4.9]$ \\
\bottomrule
\end{tabular}
\end{table}

The aggregation contrast of $+0.9$~pp, with an interval containing
zero, is not an inconclusive comparison; it is what the lemma predicts.
Both aggregations operate on a decision matrix whose effective weights
have already been fixed by the normalisation.
The distortion occurs upstream of the aggregation operator, so replacing
one operator with another cannot recover it.
The applied siting literature routinely compares aggregation methods
while leaving the normalisation implicit; on this evidence that is a
comparison of the least consequential design choice.

\begin{remark}[why pooling aggregation families obscures the mechanism]
\label{rem:pooling}
Averaging $Q_2$ across all aggregation families destroys the monotone
ordering of Table~\ref{tab:tvorder}, and does so for reasons internal to
the mechanism.
VIKOR normalises by ratios to the ideal and anti-ideal point, in which
a common scale factor cancels; VIKOR is therefore invariant to vector
normalisation, the VIKOR$\times$vector cell duplicates
VIKOR$\times$min--max, and pooling it inflates the vector average
by~$5.9$~pp.
Rank normalisation under distance-based aggregation degrades performance
far more than the additive re-weighting accounts for, because ranks
distort the Euclidean geometry on which those methods depend.
Both effects are outside the scope of Lemma~\ref{lem:amplitude} and
both are predicted by it as exclusions.
We therefore report the additive cells as the test of the mechanism and
the full factorial as a description of practice.
\end{remark}

% ----------------------------------------------------------
\subsection{Normalisation \texorpdfstring{$\times$}{x} assembly:
  optimisation amplifies the distortion}
\label{sec:interaction}
% ----------------------------------------------------------

The assembly main effect of $+16.0$~pp averages two very different
regimes.
Under min--max normalisation, replacing greedy with optimal assembly is
worth $+29.8$~pp $[+27.4,+32.2]$; under vector normalisation it is
worth $+2.2$~pp $[-0.5,+4.9]$, an interval containing zero.

The lemma explains the asymmetry.
Optimal assembly maximises $\sum_{i\in t}s(i)$ exactly, and by
Proposition~\ref{prop:p1} exact maximisation of an additive score yields
a Pareto-efficient portfolio \emph{under the weights that score
encodes}.
When the score encodes $w^{\mathrm{eff}}\neq w$, optimal assembly
solves the wrong problem exactly: it pursues the criterion trade-offs of
$w^{\mathrm{eff}}$, which in the vector case inflates \texttt{perm}
by $2.865\times$ and deflates \texttt{urban} to $0.525\times$.
Greedy assembly, being myopic, is partly insulated --- it fails for its
own reasons, but it does not optimise the distortion.

This has a direct implication for practice.
Adding integer programming on top of a ranking is a standard upgrade
\citep{BransMareschal1992,Mavrotas2006,VetscheraDeAlmeida2012},
but on this evidence the upgrade pays only when the normalisation
already preserves the elicited weights.
Exact assembly on a vector-normalised score captures almost none of the
available gain.

% ----------------------------------------------------------
\subsection{The reference set as a third instantiation of the lemma}
\label{sec:refset}
% ----------------------------------------------------------

Min--max normalisation satisfies $A_j\equiv1$ over the set on which its
constants are computed.
If that set includes alternatives that cannot be selected, the guarantee
no longer applies to the decision set: over the $\Nvia$ admissible
sites, the attainable amplitude becomes
\begin{equation}
  \label{eq:refsetA}
  A_j = \frac{\max_{\Nvia}(j)-\min_{\Nvia}(j)}
             {\max_{\Ncand}(j)-\min_{\Ncand}(j)} \;\le\; 1,
\end{equation}
which varies across criteria and therefore re-weights.
The reference-set choice is thus the same mechanism as
Sections~\ref{sec:lemmaconf}--\ref{sec:interaction}, not a separate
phenomenon.

Computing~\eqref{eq:refsetA} for the eleven criteria yields the
distortions in Table~\ref{tab:amprefset}:
five criteria are unaffected ($A_j=1.000$, because the same sites
determine the range in both sets), while \texttt{perm}~($A_j=0.482$),
\texttt{slope}~($0.683$), \texttt{rivers}~($0.648$), and
\texttt{urban}~($0.755$) lose amplitude.
The aggregate distortion is $\mathrm{TV}=0.107$, between the midranks
value ($0.091$) and the vector value ($0.312$), placing the
reference-set effect firmly within the ordering of
Table~\ref{tab:tvorder}.

\begin{table}[t]
\centering
\caption{Amplitude ratio~\eqref{eq:refsetA} by criterion (reference set
  $\Ncand\to\Nvia$). Five criteria are unaffected; four lose amplitude
  because the $\Nexcl$ excluded sites extend the range.}
\label{tab:amprefset}
\begin{tabular}{lcc}
\toprule
Criterion & $A_j$ & $w_j$ \\
\midrule
\texttt{perm}      & 0.482 & 0.152 \\
\texttt{slope}     & 0.683 & 0.113 \\
\texttt{rivers}    & 0.648 & 0.091 \\
\texttt{urban}     & 0.755 & 0.166 \\
\texttt{coast}     & 0.924 & 0.069 \\
\texttt{protected} & 0.999 & 0.066 \\
remaining 5        & 1.000 & ---   \\
\midrule
$\mathrm{TV}$      & \multicolumn{2}{c}{$0.107$} \\
\bottomrule
\end{tabular}
\end{table}

The three instantiations of the lemma now form a coherent ordering:
min--max with correct reference set ($\mathrm{TV}=0$, $Q_2=82.3\%$),
min--max over inadmissible sites ($\mathrm{TV}=0.107$, observed
$Q_2=15\text{--}65\%$ depending on assembly and comparison direction),
midranks ($\mathrm{TV}=0.091$, $Q_2=76.5\%$), and vector
($\mathrm{TV}=0.312$, $Q_2=55.5\%$).

Table~\ref{tab:fourthcell} separates the set over which the
\emph{method} computes its constants from the set over which the
\emph{evaluator} builds the reference set.

\begin{table}[t]
\centering
\caption{Reference-set misalignment, $Q_2$, SAW with min--max
  normalisation, $\kappa=200$.}
\label{tab:fourthcell}
\begin{tabular}{llcc}
\toprule
Method set & Evaluator & greedy & optimal \\
\midrule
$\Ncand$ & $A$  & 68.4\% & \textbf{100.0\%} \\
$\Ncand$ & $A'$ & 15.0\% &  20.0\% \\
$\Nvia$  & $A$  & 46.6\% &  95.0\% \\
$\Nvia$  & $A'$ & 64.6\% & \textbf{100.0\%} \\
\bottomrule
\end{tabular}
\end{table}

The pattern is a crossover: each method set performs best under the
evaluator sharing it.
Both diagonal cells reach $100\%$ under optimal assembly, as
Proposition~\ref{prop:p1} requires, which confirms the implementation
and simultaneously delimits interpretation --- the diagonal is
guaranteed, the off-diagonal is informative.

The off-diagonal costs are asymmetric.
A method normalising over all $\Ncand$ candidates but evaluated on the
admissible space loses $49.6$~pp under greedy and $80.0$~pp under
optimal assembly; the reverse misalignment costs $21.8$~pp and $5.0$~pp.
This asymmetry is consistent with~\eqref{eq:refsetA}: including the
$\Nexcl$ inadmissible sites extends the observed range, driving $A_j$
further below unity for the affected criteria and thereby increasing
$\mathrm{TV}$.
A confound is present --- objective ranges are wider in space~$A$
because it contains more candidates, contributing to the asymmetry
independently of the inadmissibility argument --- and we acknowledge it
as a limitation.

Two readings must be kept apart.
The table establishes that misalignment is costly and that the cost is
larger when inadmissible sites extend the scale, consistent with the
lemma.
It does not establish which evaluator is normatively correct.
That argument is external: it rests on the inadmissibility of the
$\Nexcl$ excluded sites (Section~\ref{sec:methods}).
\todo{D2: cite the Cape Verde legal instrument establishing the
  protected-area prohibition to convert the analytical claim into a
  normative one.}

% ----------------------------------------------------------
\subsection{The greedy shortfall}
\label{sec:greedy}
% ----------------------------------------------------------

Proposition~\ref{prop:p1} guarantees that optimal assembly of an
additive score attains the reference set under matching weights;
Table~\ref{tab:factorial} shows greedy assembly of the same score
reaching only $64.6\%$.
The shortfall is attributable to the assembly rule alone, with
aggregation, normalisation and reference set held fixed.

Draw~$302$ ($\kappa=200$) exhibits it concretely.
The greedy portfolio is $\{184,279,320\}$ with $\mathrm{gap}=0.028$:
it is dominated by an attainable portfolio by $2.8\%$ of the achievable
range on every objective simultaneously, while optimal assembly on the
identical score and weights attains the reference set exactly.

The mechanism is the separation constraint.
Greedy commits to the highest-scoring site and then discards every site
within $d_{\min}$ of it, foreclosing combinations a global search
retains.
Because~\eqref{eq:rsep} is a hard constraint rather than an objective
\citep{Erkut1996,MoonChaudhry1984}, the foreclosure cannot be undone by
re-weighting: no choice of $w$ restores a site greedy has already
excluded.
This distinguishes the assembly effect from the normalisation effect
--- the former is a combinatorial loss, the latter a re-parameterisation
of the value model.
